\chapter{विद्युत परिपथ और सुरक्षा}\label{ch:g6:circuits-and-safety}

तीन साल के फंदों, बल्बों और \omterm{def:g5:series-parallel-circuits:parallel}{शाखाओं} ने तुम्हें परिपथ बनाने वाला बना
दिया है। इस साल तुम कुछ और दुर्लभ बनते हो: परिपथ \emph{लिखने} वाले।
हर देश के बिजली वाले एक ही बिना-शब्दों की भाषा पढ़ते और लिखते हैं ---
और उसका व्याकरण, तथा मुश्किल से कमाई गई सुरक्षा-समझ का एक अध्याय, ही
चतुर तारबंदी और जली उँगलियों के बीच खड़ा रहता है।

\section{बिना शब्दों की भाषा}

\begin{definition}[परिपथ आरेख]\label{def:g6:circuits-and-safety:diagram}
\emph{परिपथ आरेख}\index{परिपथ आरेख} किसी परिपथ का मानक चित्र है: हर
उपकरण अपने माने हुए \emph{संकेत}\index{संकेत} से दिखाया जाता है, और
तार समकोण मोड़ों वाली सीधी रेखाओं से। आरेख यह दिखाता है कि परिपथ
कैसे \emph{जुड़ा} है --- कौन किसे खिलाता है, कहाँ राह फूटती है ---
और जान-बूझकर यह अनदेखा कर देता है कि असली तार मेज़ पर किस तरह
साँप-सी \omterm{def:g2:pushes-and-pulls:force}{बल} खाते हैं।
\end{definition}

\begin{omfigure}
\begin{tikzpicture}[scale=0.8]
  \draw (0.4,3.4) to[battery1] (2.4,3.4);
  \node[right, font=\small] at (2.7,3.4)
    {बैटरी (लंबी लकीर: $+$)};
  \draw (0.4,2.2) to[lamp] (2.4,2.2);
  \node[right, font=\small] at (2.7,2.2) {लैंप};
  \draw (0.4,1.0) to[spst] (2.4,1.0);
  \node[right, font=\small] at (2.7,1.0) {स्विच, खुला};
  \draw (7.4,3.4) to[cspst] (9.4,3.4);
  \node[right, font=\small] at (9.7,3.4) {स्विच, बंद};
  \draw (7.4,2.2) -- (8.05,2.2);
  \draw[thick] (8.4,2.2) circle (0.35);
  \node[font=\small] at (8.4,2.2) {M};
  \draw (8.75,2.2) -- (9.4,2.2);
  \node[right, font=\small] at (9.7,2.2) {मोटर};
  \draw (7.4,1.0) -- (8.05,1.0);
  \draw[thick] (8.4,1.0) circle (0.35);
  \draw[thick] (8.16,0.76) -- (8.64,1.24);
  \draw (8.75,1.0) -- (9.4,1.0);
  \node[right, font=\small] at (9.7,1.0) {बज़र};
\end{tikzpicture}

{\small वर्णमाला: छह \omterm{def:g6:circuits-and-safety:diagram}{संकेत} इस पुस्तक के लगभग हर परिपथ को सँभाल लेते
हैं। यही निशान दुनिया भर की कार्यशालाओं में एक जैसे पढ़े जाते हैं।}
\end{omfigure}

\begin{method}[मेज़ से काग़ज़ तक, और वापस]\label{met:g6:circuits-and-safety:translate}
बने हुए परिपथ को \emph{बनाकर लिखने} के लिए:
\begin{enumerate}
  \item उपकरणों की सूची बनाओ, और हर एक का \omterm{def:g6:circuits-and-safety:diagram}{संकेत} चुनो;
  \item \omterm{def:g3:first-electric-circuit:battery}{बैटरी} के $+$ से शुरू करके उँगली से असली फंदे का पीछा करो, और
        उपकरणों का क्रम तथा हर फूट नोट करते जाओ;
  \item इसे एक सलीक़ेदार आयत की तरह फिर से बनाओ: \omterm{def:g6:circuits-and-safety:diagram}{संकेत} दूर-दूर, तार
        सीधे, मोड़ समकोण --- जोड़ वही, साँप कोई नहीं।
\end{enumerate}
और किसी आरेख से \emph{बनाने} के लिए यही अनुवाद उलटी दिशा में करो ---
और उसके बाद, हमेशा की तरह जाँचो: $+$ से चलाया गया हर फंदा कम से कम
एक उपकरण में से होकर $-$ तक पहुँचना चाहिए।
\end{method}

\begin{example}[एक ही परिपथ, दो तसवीरें]\label{ex:g6:circuits-and-safety:portraits}
मेज़ पर: एक \omterm{def:g3:first-electric-circuit:battery}{बैटरी}, पेंसिल-डिब्बे के ऊपर से गुज़रते मगरमच्छ-क्लिप वाले
तारों की उलझन, किताब के सहारे टिका एक लैंप, और लटकता हुआ एक \omterm{ex:g3:first-electric-circuit:switch}{स्विच}।
काग़ज़ पर: एक शांत आयत --- बाईं ओर \omterm{def:g3:first-electric-circuit:battery}{बैटरी}, ऊपर \omterm{ex:g3:first-electric-circuit:switch}{स्विच}, दाईं ओर लैंप।
उलझन और आयत \emph{एक ही परिपथ} हैं: हर जोड़ हूबहू वही। दो परिपथ एक हैं
या नहीं, यह तय करने के लिए सिर्फ़ एक सवाल पूछो: कौन किससे जुड़ा है?
\end{example}

\begin{example}[मोटर और बज़र भी खेल में]\label{ex:g6:circuits-and-safety:motor}
तुम्हारे डिब्बे के दो नए साथी आरेख में ठीक लैंप जैसा ही बरताव करते
हैं: \emph{मोटर}\index{मोटर} (\omterm{def:g6:circuits-and-safety:diagram}{संकेत}: वृत्त में M) तब घूमती है जब उसमें
से धारा बहे --- पंखे, खिलौना गाड़ियाँ, ड्रिल; और
\emph{बज़र}\index{बज़र} (चोट खाया वृत्त) तब बजता है जब उसे धारा मिले
--- दरवाज़े की घंटियाँ, चेतावनी-घंटियाँ। श्रेणी और समांतर के नियम उन पर
बिना बदले लागू होते हैं: \omterm{def:g4:series-circuits:series}{श्रेणी में} लगा \omterm{ex:g3:first-electric-circuit:switch}{स्विच} उन्हें हुक्म देता है;
और समांतर शाखाएँ उन्हें एक-दूसरे से आज़ाद कर देती हैं।
\end{example}

\section{नई भाषा में पुराने नियम}

\begin{example}[माहिर की तरह आरेख पढ़ना]\label{ex:g6:circuits-and-safety:reading}
कोई भी आरेख मिले, उससे \omterm{def:g5:series-parallel-circuits:parallel}{समांतर परिपथ} वाले अध्याय की उँगली-यात्रा से
जिरह करो: $+$ से चलकर $-$ तक की हर सड़क चलो। एक ही सड़क पर बैठे सारे
उपकरण: श्रेणी --- एक दरार सबको रोक देती है, क्रम आज़ाद है, चमक बँटी
हुई। \omterm{def:g5:series-parallel-circuits:parallel}{संधि} के बिंदुओं पर फूटती सड़कें: समांतर --- शाखाएँ आज़ाद, हर एक
पूरी चमक पर, और अब \omterm{ex:g3:first-electric-circuit:switch}{स्विच} की जगह सचमुच का फ़ैसला। भाषा बदली; नियम
नहीं।
\end{example}

\begin{omfigure}
\begin{tikzpicture}[scale=0.8]
  \draw (0,0) to[battery1] (0,3.2)
    to[cspst] (2.6,3.2);
  \draw (2.6,3.2) to[short] (2.6,2.4)
    to[lamp] (5.8,2.4) to[short] (5.8,3.2);
  \draw (2.6,3.2) to[short] (2.6,4.0) -- (3.4,4.0);
  \draw[thick] (3.75,4.0) circle (0.35);
  \node[font=\small] at (3.75,4.0) {M};
  \draw (4.1,4.0) -- (5.8,4.0) to[short] (5.8,3.2);
  \draw (5.8,3.2) to[short] (7.4,3.2) to[short] (7.4,0)
    to[short] (0,0);
  \fill (2.6,3.2) circle (2pt);
  \fill (5.8,3.2) circle (2pt);
\end{tikzpicture}

{\small बनाने से पहले पढ़ो: साझा सड़क पर एक \omterm{ex:g3:first-electric-circuit:switch}{स्विच}, फिर एक फूट ---
एक \omterm{def:g5:series-parallel-circuits:parallel}{शाखा} पर लैंप, दूसरी पर पंखे की \omterm{ex:g6:circuits-and-safety:motor}{मोटर}। \omterm{ex:g3:first-electric-circuit:switch}{स्विच} दोनों को हुक्म देता है;
और शाखाएँ एक-दूसरे की मुसीबतों से बेपरवाह रहती हैं।}
\end{omfigure}

\section{निशानों वाला अध्याय}

बिजली के सुरक्षा-नियमों का दाम आगों और उनसे भी बुरी चीज़ों से चुकाया
गया है। वे नियम यहाँ हैं, अपनी वजहों के साथ --- पेशेवर दोनों जानता है।

\begin{definition}[लघुपथन]\label{def:g6:circuits-and-safety:short}
\emph{लघुपथन}\index{लघुपथन} तब होता है जब कोई नंगी \omterm{def:g4:conductors-insulators:conductor}{चालक} सड़क \omterm{def:g3:first-electric-circuit:battery}{बैटरी} के
(या सॉकेट के) दोनों पक्षों को \emph{सीधे} जोड़ देती है और हर उपकरण को
छोड़ देती है। रास्ते में कोई काम की चीज़ न हो तो धारा बेतहाशा दौड़
पड़ती है --- तार तपते हैं, \omterm{def:g4:conductors-insulators:insulator}{विद्युतरोधी} परत \omterm{def:g2:ice-water-steam:meltfreeze}{पिघलती} है, बैटरियाँ बरबाद
होती हैं, और घर की तारबंदी में आग तक लग सकती है। अगले साल का परिपथ वाला
अध्याय इस खलनायक को पास से देखेगा; इस साल इसका चेहरा पहचान लो और इसे
अपने परिपथों से दूर रखो।
\end{definition}

\begin{remark}[घर के अंगरक्षक]\label{rem:g6:circuits-and-safety:guards}
घर की फ़्यूज़-पेटी में झाँको: छोटे-छोटे स्विचों की क़तारें ---
\emph{परिपथ विच्छेदक} --- घर के विशाल \omterm{def:g5:series-parallel-circuits:parallel}{समांतर परिपथ} की हर \omterm{def:g5:series-parallel-circuits:parallel}{शाखा} के लिए
एक। हर एक अपनी \omterm{def:g5:series-parallel-circuits:parallel}{शाखा} पर पहरा देता है: अगर किसी ख़राबी से धारा दौड़
पड़े --- \omterm{def:g6:circuits-and-safety:short}{लघुपथन}, या हद से ज़्यादा लदा सॉकेट --- तो पहरेदार पलक झपकते
ख़ुद को बंद कर लेता है और तार तपने से पहले \omterm{def:g5:series-parallel-circuits:parallel}{शाखा} काट देता है। जब कोई
विच्छेदक ``गिर जाता'' है, तो वह बदमाशी नहीं कर रहा: उसने अभी-अभी अपना
इकलौता काम किया है। उसे दोबारा चढ़ाने से पहले ख़राबी ढूँढ़ो।
\end{remark}

\begin{remark}[नियम, वजहों सहित]\label{rem:g6:circuits-and-safety:rules}
\begin{enumerate}
  \item \emph{प्रयोग सिर्फ़ \omterm{def:g3:first-electric-circuit:battery}{बैटरी} से; सॉकेट से कभी नहीं।} घर की
        बिजली सैकड़ों गुना ज़्यादा ज़ोर से धकेलती है --- नम त्वचा के
        रास्ते वह दिल तक रोक सकती है।
  \item \emph{\omterm{def:g3:first-electric-circuit:battery}{बैटरी} के दोनों सिरों को सीधे जोड़ता कोई नंगा तार नहीं।} वह \omterm{def:g6:circuits-and-safety:short}{लघुपथन} है:
        बरबाद सेल, तपता तार, जली उँगलियाँ।
  \item \emph{पानी और घर की बिजली कभी एक-दूसरे की पहुँच में न आएँ:}
        नहाने की जगह के पास कोई यंत्र नहीं, गीले हाथों से कोई प्लग
        नहीं --- नम त्वचा \omterm{def:g4:conductors-insulators:conductor}{चालक} है।
  \item \emph{सॉकेट में कभी कुछ मत घुसाओ।} फ़्यूज़-पेटी का पहरेदार
        तेज़ है, पर उस पर दाँव कभी मत लगाओ।
  \item \emph{ख़राब तार कूड़ेदान में जाते हैं,} दराज़ में नहीं: चटकी
        हुई \omterm{def:g4:conductors-insulators:insulator}{विद्युतरोधी} परत दरार वाली बाड़ है।
\end{enumerate}
\end{remark}

\section{अभ्यास}

\begin{exercise}[$\star$]\label{exo:g6:circuits-and-safety:1}
\omterm{def:g6:circuits-and-safety:diagram}{परिपथ आरेख} किस बात को ईमानदारी से दिखाता है, और किसे जान-बूझकर अनदेखा
कर देता है?
\end{exercise}

\begin{exercise}[$\star$]\label{exo:g6:circuits-and-safety:2}
याद से \omterm{def:g6:circuits-and-safety:diagram}{संकेत} बनाओ: \omterm{def:g3:first-electric-circuit:battery}{बैटरी}; \ लैंप; \ खुला \omterm{ex:g3:first-electric-circuit:switch}{स्विच}; \ बंद \omterm{ex:g3:first-electric-circuit:switch}{स्विच}; \
\omterm{ex:g6:circuits-and-safety:motor}{मोटर}; \ \omterm{ex:g6:circuits-and-safety:motor}{बज़र}। \omterm{def:g3:first-electric-circuit:battery}{बैटरी} का $+$ पक्ष कौन-सा निशान बताता है?
\end{exercise}

\begin{exercise}[$\star$]\label{exo:g6:circuits-and-safety:3}
दरवाज़े की घंटी का आरेख बनाओ: एक ही फंदे में \omterm{def:g3:first-electric-circuit:battery}{बैटरी}, दबाने वाला \omterm{ex:g3:first-electric-circuit:switch}{स्विच}
और \omterm{ex:g6:circuits-and-safety:motor}{बज़र}। उँगली हटते ही घंटी क्यों रुक जाती है?
\end{exercise}

\begin{exercise}[$\star$]\label{exo:g6:circuits-and-safety:4}
एक ही डिब्बे के पुरज़ों से दो परिपथ बनाए गए हैं: क्या वे ``एक ही
परिपथ'' हैं? कौन-सा अकेला सवाल फ़ैसला कर देता है?
\end{exercise}

\begin{exercise}[$\star$]\label{exo:g6:circuits-and-safety:5}
अध्याय वाले लैंप-और-मोटर आरेख में लैंप फुँक जाए तो पंखे का क्या होगा?
और \omterm{ex:g3:first-electric-circuit:switch}{स्विच} खुल जाए तो दोनों का? कौन-से पुराने नियम उत्तर देते हैं?
\end{exercise}

\begin{exercise}[$\star$]\label{exo:g6:circuits-and-safety:6}
\omterm{def:g6:circuits-and-safety:short}{लघुपथन} क्या है, और उससे तार क्यों तपते हैं? उसके दो शिकार बताओ।
\end{exercise}

\begin{exercise}[$\star$]\label{exo:g6:circuits-and-safety:7}
फ़्यूज़-पेटी का एक परिपथ विच्छेदक गिर गया है। उसने अभी-अभी घर के लिए
क्या किया? उसे दोबारा चढ़ाने से पहले क्या करना चाहिए?
\end{exercise}

\begin{exercise}[$\star$]\label{exo:g6:circuits-and-safety:8}
हर नियम के पीछे की वजह बताओ: नहाने की जगह के पास कोई यंत्र नहीं; \
ख़राब तार कूड़ेदान में; \ \omterm{def:g3:first-electric-circuit:battery}{बैटरी} के खेल हाँ, सॉकेट कभी नहीं।
\end{exercise}

\begin{exercise}[$\star\star$]\label{exo:g6:circuits-and-safety:9}
इसे आरेख में बदलो (शब्दों में या चित्र बनाकर): एक \omterm{def:g3:first-electric-circuit:battery}{बैटरी}; उसके $+$ से
एक तार बंद \omterm{ex:g3:first-electric-circuit:switch}{स्विच} तक; \omterm{ex:g3:first-electric-circuit:switch}{स्विच} से एक फूट --- एक \omterm{def:g5:series-parallel-circuits:parallel}{शाखा} \omterm{ex:g6:circuits-and-safety:motor}{बज़र} से होकर, दूसरी
लैंप से होकर; शाखाएँ फिर मिलकर $-$ तक लौटती हैं। बताओ कि क्या सुनाई
और क्या दिखाई देगा, और अगर सिर्फ़ \omterm{ex:g6:circuits-and-safety:motor}{बज़र} वाली \omterm{def:g5:series-parallel-circuits:parallel}{शाखा} को उसका अपना एक और
\omterm{ex:g3:first-electric-circuit:switch}{स्विच} मिल जाए और वह खुला हो, तो क्या बदलेगा।
\end{exercise}

\begin{exercise}[$\star\star$]\label{exo:g6:circuits-and-safety:10}
एक खिलौना पंखा इस तरह जोड़ा गया है: \omterm{def:g3:first-electric-circuit:battery}{बैटरी} का $+$ \omterm{ex:g6:circuits-and-safety:motor}{मोटर} तक, \omterm{ex:g6:circuits-and-safety:motor}{मोटर} से
\omterm{def:g3:first-electric-circuit:battery}{बैटरी} का $-$, और --- ``मज़बूती के लिए'' --- $+$ से $-$ तक सीधे एक
अतिरिक्त नंगा तार। पंखा बमुश्किल घूमता है और वह अतिरिक्त तार गरम होने
लगता है। \cref{def:g6:circuits-and-safety:short} से जाँच करो: धारा
कौन-सी सड़क पसंद कर रही है, और यह जोड़ ख़तरनाक क्यों है?
\end{exercise}

\begin{exercise}[$\star\star$]\label{exo:g6:circuits-and-safety:11}
दादाजी की कार्यशाला वाली आदत: बत्ती के किसी भी फ़िटिंग को छूने से
पहले वे उसका दीवार वाला \omterm{ex:g3:first-electric-circuit:switch}{स्विच} बंद करते हैं \emph{और} फ़्यूज़-पेटी में
उस पूरी \omterm{def:g5:series-parallel-circuits:parallel}{शाखा} का विच्छेदक भी। श्रेणी वाली सोच से समझाओ कि हर कटाव क्या
हासिल करता है --- और सावधान आदमी दोनों क्यों करता है।
\end{exercise}

\begin{exercise}[$\star\star\star$]\label{exo:g6:circuits-and-safety:12}
कठपुतली रंगमंच की तारबंदी बनाओ: एक मुख्य \omterm{ex:g3:first-electric-circuit:switch}{स्विच} जो सब कुछ अँधेरा कर दे;
मंच के दो लैंप जो एक-दूसरे से आज़ाद रहकर ख़राब हो सकें; और परदे की एक
\omterm{ex:g6:circuits-and-safety:motor}{मोटर} तथा शो शुरू होने का एक \omterm{ex:g6:circuits-and-safety:motor}{बज़र}, हर एक अपने अलग नियंत्रण-स्विच के
साथ। आरेख बताओ (या बनाओ) --- कौन-से उपकरण साझा सड़क पर सवार हैं,
कौन-से \omterm{def:g5:series-parallel-circuits:parallel}{शाखाओं} पर, और हर \omterm{ex:g3:first-electric-circuit:switch}{स्विच} कहाँ बैठता है --- और अपने डिज़ाइन को
श्रेणी तथा समांतर के हर उस नियम पर परखो जो तुम्हारे पास है।
\end{exercise}

\section{समस्या: बग़ीचे की कोठरी में तारबंदी}

\begin{problem}[{सप्ताहांत समस्या --- एक परिवार अपनी बग़ीचे वाली कोठरी
में तार बिछाता है; पेचकस से पहले आरेख; और निरीक्षक की
सूची}]\label{pb:g6:circuits-and-safety:1}
परिवार कोठरी के \omterm{def:g3:first-electric-circuit:battery}{बैटरी} से चलने वाले छोटे जाल की योजना पहले काग़ज़ पर
बनाता है --- जैसा पेशेवर लोग करते हैं।

\medskip\noindent\textbf{भाग I --- काग़ज़ पर योजना।}
इच्छा-सूची: दरवाज़े से चलने वाला छत का एक लैंप; अपने अलग \omterm{ex:g3:first-electric-circuit:switch}{स्विच} वाला
कार्यमेज़ का एक लैंप; और अपने अलग \omterm{ex:g3:first-electric-circuit:switch}{स्विच} वाला \omterm{def:g2:air-around-us:air}{हवा} निकालने का छोटा पंखा
(\omterm{ex:g6:circuits-and-safety:motor}{मोटर})।
\begin{enumerate}
  \item तीनों उपकरण एक श्रेणी फंदे के बजाय समांतर \omterm{def:g5:series-parallel-circuits:parallel}{शाखाओं} पर क्यों
        होने चाहिए? ऐसी दो सुविधाएँ बताओ जो \omterm{def:g4:series-circuits:series}{श्रेणी में} टूट जातीं।
  \item हर \omterm{def:g5:series-parallel-circuits:parallel}{शाखा} पर उसका अपना \omterm{ex:g3:first-electric-circuit:switch}{स्विच} है। हर \omterm{ex:g3:first-electric-circuit:switch}{स्विच} किसे हुक्म देता है,
        और किस चीज़ पर इनमें से किसी का हुक्म नहीं चलता?
  \item पूरी कोठरी को एक साथ अँधेरा और चुप कर देने वाला मुख्य \omterm{ex:g3:first-electric-circuit:switch}{स्विच}
        कहाँ बैठना चाहिए?
  \item पूरा आरेख बनाओ (या ठीक-ठीक बताओ): \omterm{def:g3:first-electric-circuit:battery}{बैटरी}, मुख्य \omterm{ex:g3:first-electric-circuit:switch}{स्विच}, तीन
        शाखाएँ, और कुल चार \omterm{ex:g3:first-electric-circuit:switch}{स्विच}।
\end{enumerate}

\medskip\noindent\textbf{भाग II --- बनकर तैयार, और जाँच में।}
\begin{enumerate}[resume]
  \item बनाकर चालू करने पर छत का लैंप जलता तो है --- पर मंद, और
        कार्यमेज़ का लैंप सिर्फ़ तब जलता है जब \emph{पंखे} का \omterm{ex:g3:first-electric-circuit:switch}{स्विच}
        बंद हो। उँगली-यात्रा से पता चलता है कि मेज़ वाला लैंप और
        पंखा एक ही \omterm{def:g5:series-parallel-circuits:parallel}{शाखा} पर, एक के पीछे एक बैठे हैं। उनकी तरतीब का
        नाम बताओ, और दोनों लक्षण पुराने नियमों से समझाओ।
  \item योजना बहाल करने वाली एक-तार वाली मरम्मत बताओ।
  \item सफ़ाई करते समय एक पेचकस पल भर के लिए \omterm{def:g3:first-electric-circuit:battery}{बैटरी} के दोनों नंगे
        सिरों पर पड़ा रहता है और गरम हो जाता है। इस घटना का नाम
        बताओ, और यह भी कि परिवार को उसके बाद \omterm{def:g3:first-electric-circuit:battery}{बैटरी} के बारे में क्या
        जाँचना चाहिए।
  \item निरीक्षक बने अभिभावक एक नियम जोड़ते हैं: ``हर जोड़ लपेटा हुआ,
        कहीं भी नंगा ताँबा नहीं।'' यह अकेला नियम किन दो ख़तरों से
        बचा रहा है?
\end{enumerate}

\medskip\noindent\textbf{भाग III --- तूफ़ान वाली रात।}
\begin{enumerate}[resume]
  \item बरसाती रात: कोठरी की छत से कार्यमेज़ पर पानी टपकता है।
        \omterm{def:g3:first-electric-circuit:battery}{बैटरी} से चलने वाली कोठरी में भी परिवार को गीले हाथों से
        तारबंदी क्यों नहीं छूनी चाहिए --- और घर की मुख्य बिजली वाली
        तारबंदी में वही स्पर्श कहीं ज़्यादा भारी क्यों पड़ता?
  \item घर की बत्तियाँ टिमटिमाती हैं और एक विच्छेदक गिर जाता है।
        विच्छेदक ने क्या पकड़ा, और कामों का समझदार क्रम क्या है?
  \item सुबह: परिवार दरवाज़े की घंटी जोड़ने पर बहस करता है --- कोठरी
        पर \omterm{ex:g6:circuits-and-safety:motor}{बज़र}, घर पर दबाने वाला \omterm{ex:g3:first-electric-circuit:switch}{स्विच}, और बीच में लंबा दुहरा तार।
        कोठरी के जाल के साथ \omterm{def:g4:series-circuits:series}{श्रेणी में} या \omterm{def:g5:series-parallel-circuits:parallel}{समांतर में}? उसकी धारा कहाँ
        से आएगी? फंदे का रेखाचित्र बनाओ।
  \item कोठरी के लिए परिवार का तीन पंक्तियों वाला सुरक्षा-घोषणापत्र
        लिखो, हर पंक्ति में एक बात: सॉकेट, पानी, नंगे तार।
\end{enumerate}
\end{problem}
